% Standalone Data Availability Statement (PCCP / RSC)
% Compile: pdflatex data_availability.tex
\documentclass[11pt,a4paper]{article}
\usepackage[margin=2.5cm]{geometry}
\usepackage{times}
\usepackage{hyperref}
\usepackage{xcolor}
\hypersetup{colorlinks=true, urlcolor=blue, citecolor=blue}
\setlength{\parskip}{6pt}
\setlength{\parindent}{0pt}

\begin{document}

\begin{center}
{\Large\bfseries Data Availability Statement}\\[0.6em]
{\small For the manuscript: ``What a radical-pair quantum reservoir would require:
readable capacity, no quantum advantage, and a nuclear-register criterion''}\\[0.3em]
{\small Hikaru Wakaura and Taiki Tanimae, QIRI (Quantum Integrated Research Institute Inc.), Tokyo, Japan}\\[0.3em]
{\small \texttt{h.wakaura@qiri.co.jp}, \texttt{t.tanimae@qiri.co.jp}}
\end{center}

\vspace{0.8em}
\hrule
\vspace{1em}

\noindent\textbf{Code.}
All results in this work are produced by the open-source Python package
\texttt{qbscreen} (MIT licence), available at
\url{https://github.com/deeptell-inc/new_Q_brain}. The package
implements the exact Liouville-space Haberkorn master equation with electronic
decoherence and the correlated-pair-birth quantum-reservoir-computing protocol. The
specific modules are:

\begin{itemize}\setlength{\itemsep}{2pt}
\item \texttt{qbscreen/master\_equation.py} --- Liouville-space solver and its
validation against the closed-form coherent limit (ESI, S1);
\item \texttt{qbscreen/reservoir.py} --- reservoir Hamiltonian and propagator, and
the out-of-sample information-processing capacity (IPC) and memory capacity (MC)
estimators (trained on the first half of each run, scored on the held-out second
half);
\item \texttt{qbscreen/corrected\_injection.py} --- the chemically consistent
correlated-pair birth channel $\rho_{\mathrm e}(s)=s\,|S\rangle\langle S|
+(1-s)P_{\mathrm T}/3$ and its S--$\mathrm{T}_0$ coherent variant, together with the
single-electron-reset control that reproduces the artifact it replaces;
\item \texttt{qbscreen/readout\_routes.py} --- the seven biologically accessible
readout routes with recombination active, the product-weighted CIDNP observable, the
turnover clock, and the nuclear-register depolarisation parameter $q$;
\item \texttt{qbscreen/final\_numbers.py} --- protocol-consistent recomputation of
every headline quantity reported in the main text and ESI;
\item \texttt{qbscreen/qrc\_benchmarks.py} --- classical echo-state-network
baselines, the NARMA2 benchmark, and the coherence-time scan;
\item \texttt{manuscript/make\_schematic.py} and
\texttt{manuscript/make\_fig\_criterion.py} --- Figs.~1--3.
\end{itemize}

\noindent\textbf{Data.}
All numerical outputs underlying the figures and tables are provided as JSON files in
\texttt{simulation\_results/} of the repository, principally
\texttt{final\_numbers.json}, \texttt{corrected\_injection.json},
\texttt{readout\_routes.json}, \texttt{register\_reuse.json},
\texttt{grid\_convergence.json}, \texttt{qrc\_benchmarks.json} and
\texttt{cryptochrome\_reality.json}. No experimental data were generated; the study
is entirely computational. Every reported capacity is out of sample, is quoted with
its shuffled-input null floor, and is shown to be converged in sample length; values
are means over $4$--$10$ input realisations with the stated standard deviations;
the turnover-clock scan uses a single realisation, since the quantity of interest
there is the ratio of memory capacities across pause lengths on the same input.

\noindent\textbf{Reproducibility.}
The central claims---the agreement of the Liouville-space solver with the coherent
limit, the Dambre capacity bound, the correlated-pair-birth readout capacities, and
the nuclear-register reuse criterion---are covered by an automated test suite
(\texttt{qbscreen/tests/}; 28 tests, all passing). All figures and numbers can be
regenerated from a clean checkout with
\begin{quote}\small
\texttt{pip install -e .}\\
\texttt{python -m pytest qbscreen/tests/ -q}\\
\texttt{python -m qbscreen.final\_numbers}\\
\texttt{python -m qbscreen.readout\_routes}
\end{quote}

\noindent\textbf{Software.}
Python~3, NumPy ($\ge1.24$), SciPy ($\ge1.10$) and Matplotlib ($\ge3.7$). No
proprietary software is required.

\vspace{1.5em}
\hrule
\vspace{1em}

\noindent\textit{This statement follows the Royal Society of Chemistry Data Sharing
Policy: \url{https://www.rsc.org/journals-books-databases/about-journals/data-sharing-policy/}}

\end{document}
